\documentclass[12pt, a4paper]{article}
\usepackage[T1]{fontenc}
\usepackage[utf8]{inputenc}
\usepackage[spanish]{babel}
\usepackage{geometry}
\usepackage{graphicx}
\usepackage{xcolor}
\usepackage{listings}
\usepackage{booktabs}
\usepackage{hyperref}
\usepackage{palatino} % Fuente más amigable para lectura
\usepackage{textcomp} % Para símbolos adicionales como \textasciiexclamdown

\geometry{top=2.5cm, bottom=2.5cm, left=2.5cm, right=2.5cm}

% Configuración de colores y estilo para el código
\definecolor{codegreen}{rgb}{0,0.6,0}
\definecolor{codegray}{rgb}{0.5,0.5,0.5}
\definecolor{codepurple}{rgb}{0.58,0,0.82}
\definecolor{backcolour}{rgb}{0.95,0.95,0.92}

\lstdefinestyle{mystyle}{
    backgroundcolor=\color{backcolour},
    commentstyle=\color{codegreen},
    keywordstyle=\color{magenta},
    numberstyle=\tiny\color{codegray},
    stringstyle=\color{codepurple},
    basicstyle=\ttfamily\footnotesize,
    breakatwhitespace=false,
    breaklines=true,
    captionpos=b,
    keepspaces=true,
    numbers=left,
    numbersep=5pt,
    showspaces=false,
    showstringspaces=false,
    showtabs=false,
    tabsize=2
}

\lstset{style=mystyle,
    literate={á}{{\'a}}1
             {é}{{\'e}}1
             {í}{{\'i}}1
             {ó}{{\'o}}1
             {ú}{{\'u}}1
             {ñ}{{\~n}}1
             {Á}{{\'A}}1
             {É}{{\'E}}1
             {Í}{{\'I}}1
             {Ó}{{\'O}}1
             {Ú}{{\'U}}1
             {Ñ}{{\~N}}1
             {¡}{{\textexclamdown}}1 % Solución: \textexclamdown (de T1) funciona en lstlisting
}

\title{\textbf{Programación de Agentes de IA}\\ \large La Arquitectura de 3 Capas explicada con una Cocina}
\author{Clase de Informática}
\date{\today}

\begin{document}

\maketitle

\section{Introducción: El Problema del Chef Distraído}

Imaginen que están construyendo un robot para cocinar en un restaurante de lujo. El cerebro de este robot es una Inteligencia Artificial (como Gemini o GPT).

El problema es que las IAs son muy creativas, pero a veces se distraen o 'alucinan'. Son como un \textbf{Chef Genio pero distraído}. Si le pides que corte una cebolla, quizás te escriba un poema sobre las capas de la cebolla en lugar de cortarla.

Para solucionar esto y hacer que el robot sea confiable, usamos una técnica llamada \textbf{Arquitectura de 3 Capas}.

\section{La Analogía de la Cocina}

\subsection{Capa 1: Directivas (El Libro de Recetas)}
\begin{itemize}
    \item \textbf{¿Qué es?} Archivos de texto con instrucciones fijas (YAML).
    \item \textbf{Función:} Es el manual que dice \textit{qué} hacer paso a paso. Aquí no hay magia, solo reglas.
    \item \textbf{Ejemplo:} 'Paso 1: Cortar cebolla. Paso 2: Sofreír'.
\end{itemize}

\subsection{Capa 2: Orquestación (El Chef Ejecutivo / La IA)}
\begin{itemize}
    \item \textbf{¿Qué es?} El modelo de Inteligencia Artificial (el cerebro).
    \item \textbf{Función:} Es el jefe. Su trabajo \textbf{no es picar la cebolla} con sus propias manos. Su trabajo es:
    \begin{enumerate}
        \item Leer la receta (Capa 1).
        \item Dar la orden a los cocineros: '¡Necesito picar cebolla ahora!'.
        \item Verificar que el plato esté bien hecho.
        \item Si hay un error (ej. se acabaron las cebollas), decidir qué hacer.
    \end{enumerate}
\end{itemize}

\subsection{Capa 3: Ejecución (Los Cocineros de Línea / Herramientas)}
\begin{itemize}
    \item \textbf{¿Qué es?} Scripts de programación (código Python).
    \item \textbf{Función:} Son las manos y las herramientas. Son 'tontos' pero \textbf{extremadamente precisos}.
    \item \textbf{Diferencia clave:} Si la IA (Capa 2) dice 'Calcula la raíz cuadrada de 144', el Script (Capa 3) saca la calculadora y devuelve '12'. Nunca dirá 'creo que es 13'.
\end{itemize}

\section{Ejemplo Práctico en Código}

Si queremos que el agente 'Guarde un recuerdo', así se ve la Capa 3 (Ejecución) en Python. Es código simple que hace una sola cosa bien:

\begin{lstlisting}[language=Python, caption=Script de Ejecución (Capa 3)]
import sys

# Recibimos el texto del Chef (Capa 2)
texto = sys.argv[1]

# Hacemos el trabajo sucio: Guardar en disco
archivo = open('memoria.txt', 'a')
archivo.write(texto + '\n')
archivo.close()

print('¡Listo jefe! Guardado en el disco.')
\end{lstlisting}

\section{Resumen Visual}

\begin{table}[h]
\centering
\begin{tabular}{@{}llll@{}}
\toprule
\textbf{Capa} & \textbf{Nombre} & \textbf{Rol en la Cocina} & \textbf{Herramienta Real} \\ \midrule
1 & Directivas & El Libro de Recetas (El QUÉ) & Archivos YAML/Texto \\
2 & Orquestación & El Chef Ejecutivo (El CUÁNDO) & IA (Gemini/GPT) \\
3 & Ejecución & Cocineros de Línea (El CÓMO) & Python (Scripts) \\ \bottomrule
\end{tabular}
\caption{Comparativa de la Arquitectura de 3 Capas}
\end{table}

\section{Conclusión}
Separamos el \textbf{Pensamiento Probabilístico} (IA) de la \textbf{Acción Determinista} (Código). Así obtenemos lo mejor de los dos mundos: un robot que entiende lo que quieres y que nunca falla al hacerlo.

\end{document}